% !PRINTED 255
\weberpage{255}
%<<CONT
rationally in terms of those of the corresponding equation for degree
$p^2$, and substitute this expression into the transformation equation
belonging to $p$.

According to 3., formula (7), the equation
\weq{15}{x^4 + 18x^2 + \gamma_3(\omega)x - 27 = 0}
has the following two functions as roots:
\weq{16}{-\left(\frac{\eta\left(\frac{\omega}{3}\right)}{\eta(\omega)}\right)^6, \quad 27\left(\frac{\eta(3\omega)}{\eta(\omega)}\right)^6}
Let the first of these values be $x$. Under the substitution
$\left(\omega,\frac{\omega}{3}\right)$, the second becomes $\dfrac{-27}{x}$;
consequently $x$ also satisfies the equation
\weq{17}{27^3 + 18\cdot27x^2 - \gamma_3\left(\frac{\omega}{3}\right)x^3 - x^4 = 0}
If we now set
\weq{18}{y = \left(\frac{\eta\left(\frac{\omega}{9}\right)}{\eta(\omega)}\right)^3,}
then, under the substitution $\left(\omega,\frac{\omega}{3}\right)$,
$x$ becomes $\dfrac{y^2}{x}$, so that one also obtains the equation:
\weq{19}{y^8 + 18y^4x^2 + \gamma_3\left(\frac{\omega}{3}\right)y^2x^3 - 27x^4 = 0,}
and by elimination of $\gamma_3\left(\frac{\omega}{3}\right)$ from
(17) and (19) one obtains, after cancellation of the factor $y^2+27$,
\weq{20}{x^4 - 18x^2y^2 - y^2(y^4 - 27y^2 + 27^2) = 0.}

If this quadratic equation is solved for $x^2$, it follows that:
\weq{21}{x^2 = y^3 + 9y^2 + 27y = (y+3)^3 - 27,}
where the sign is decided by substituting the initial terms of the
developments, most simply, no doubt, since this equation (by § 69) is
also satisfied for
\[
x = 27\left(\frac{\eta(3\omega)}{\eta(\omega)}\right)^6, \quad y = 27\left(\frac{\eta(9\omega)}{\eta(\omega)}\right)^3
\]
when one sets
\[
x = 27q + \cdots, \quad y = 27q^2 + \cdots
\]
in it.
